\section{Mass}

Why is the top quark so much heavier than the electron. A particle's mass is set by how deep it sits in a warped
throat of the mesh. Each step deeper multiplies the scale by the warp factor, about eighteen, so a few steps of
depth produce the huge spread of masses we see, without any single absurdly fine-tuned number. The famous hierarchy
problem, a ratio of ten thousand trillion between the Planck and electroweak scales, shrinks to a small depth count of
about thirteen.

But which particle sits at which depth is not fixed by the geometry. Those depths are free, and they are exactly the
same free knobs, the Yukawa couplings, that the Standard Model also leaves free. This is not a failure. The model
explains the pattern and the shape of the hierarchy, and hands you the geometric origin of each free number, a depth,
while the absolute values stay open, just as in every theory we have. A direct test of whether the deterministic growth
secretly fixes the depths came back negative, so the freedom is genuine, a true boundary condition, not a hidden gap.

The deepest free parameter, the overall scale, is partly the wrong question. Every physical observable is a
dimensionless ratio, and an absolute scale in any unit is a choice of ruler, not a measured fact, so there is no
physical free dimensionful parameter, only a convention. The genuine free parameters are all dimensionless, and they
are the roughly nineteen the Standard Model already has, each now tied to a specific geometric origin, a depth or a
phase. The model does not remove them, it shows where each one lives.
